\begin{tikzpicture}[font=\scriptsize, >=latex]
% load/store addressing
\draw[thick, fill=accentsoft] (0.4,3.2) rectangle (2.0,4.4);
\node[align=center] at (1.2,4.0) {rs1};
\node[align=center, font=\tiny] at (1.2,3.5) {base = 0x100};
\node[font=\Large] at (2.6,3.8) {$+$};
\draw[thick, fill=graybg] (3.0,3.2) rectangle (4.2,4.4);
\node[align=center] at (3.6,4.0) {imm[11:0]};
\node[align=center, font=\tiny] at (3.6,3.5) {offset = 4};
\draw[->, very thick] (4.6,3.8) -- (5.6,3.8);
\draw[thick, fill=goodbg] (6.0,3.2) rectangle (9.2,4.4);
\node[align=center] at (7.6,4.0) {address 0x104};
\node[align=center, font=\tiny] at (7.6,3.5) {reads/writes word a[1]};
\node[align=left, text=slate, text width=10cm] at (5.0,2.3)
  {\ttfamily lw t4, 4(t0) \rmfamily is the whole addressing story:
   register $+$ sign-extended 12-bit immediate. No indexing modes,
   no scaling --- RV32 keeps the ALU cheap and lets the pointer-bump
   instruction do the walking (that's the \ttfamily addi t0,t0,4
   \rmfamily in the loop).};
\end{tikzpicture}
